% s03-4-2-tempered.tex —— 3.4.3 与 3.4.4：Iwahori–Hecke 温性与 IFLE 逐点等价（印刷页 111–125）

\subsection{Iwahori–Hecke 温性}
\label{sec:ih-tempered}

本小节引入 Iwahori–Hecke 温性这一概念，并说明如何用它证明命题~\ref{prop:3.4.2.1}。

考虑 $\Aff^{\spec}_{\check G}$ 在 $x_0$ 处的纤维：
\[
\IndCoh^{*}\Bigl(\check{\mathfrak{n}}/\check B^{\Ran}_{x_0}
\times_{(\LS^{\mer}_{\check G})_{\Ran_{x_0}}}
\check{\mathfrak{n}}/\check B^{\Ran}_{x_0}\Bigr)^{x_0}
\simeq \IndCoh\bigl(\check{\mathfrak{n}}/\check B
\times_{\check{\mathfrak{g}}/\check G}\check{\mathfrak{n}}/\check B\bigr).
\]
把范畴
\[
\QCoh\bigl(\check{\mathfrak{n}}/\check B
\times_{\check{\mathfrak{g}}/\check G}\check{\mathfrak{n}}/\check B\bigr)
\]
记作 $\Aff^{\spec,\temp}_{\check G,x_0}$。有全忠实嵌入
\[
\Aff^{\spec,\temp}_{\check G,x_0}\longrightarrow \Aff^{\spec}_{\check G,x_0},
\]
它有幺半的连续右伴随
\[
\Aff^{\spec}_{\check G,x_0}\longrightarrow \Aff^{\spec,\temp}_{\check G,x_0}.
\]
设 $M$ 是 $\Aff^{\spec}_{G,x_0}$-模范畴，则可考虑
\[
M^{\temp}=\Aff^{\spec,\temp}_{\check G,x_0}
\otimes_{\Aff^{\spec}_{\check G,x_0}}M,
\]
而伴随对
$\Aff^{\spec,\temp}_{\check G,x_0}\rightleftharpoons \Aff^{\spec}_{\check G,x_0}$
诱导出伴随对 $M^{\temp}\rightleftharpoons M$。

\begin{definition}\label{def:ih-tempered}
称 $\Aff^{\spec}_{G,x_0}$-模范畴 $M$ 是 \emph{Iwahori–Hecke 温的}，若上述伴随是等价。
\end{definition}

\begin{remark}\label{rem:ih-tempered}
我们用「Iwahori–Hecke」这一修饰语，是为了把此概念与 \cite[第 7.1 节]{ABC+24a} 中定义的
类似概念区分开。后者我们称为\emph{球面温性}。
\end{remark}
\srcpage{112}

现设
\[
\LS^{\mathrm{nilp\text{-}mon}}_{\check G,x_0}
\]
是
\[
\check{\mathfrak{n}}/\check B\longrightarrow \LS^{\mer}_{\check G,x_0}
\]
的像，注意
\[
\Aff^{\spec,\temp}_{\check G,x_0}
\simeq
\Fun_{\QCoh((\LS^{\mer}_{\check G,x_0})^{\wedge\mathrm{nilp\text{-}mon}})}
\bigl(\QCoh((\check{\mathfrak{n}}/\check B)_{x_0}),
\QCoh((\check{\mathfrak{n}}/\check B)_{x_0})\bigr).
\]
对 $\Aff^{\spec}_{G,x_0}$-模范畴 $M$，由于
$\QCoh((\check{\mathfrak{n}}/\check B)_{x_0})$ 是刚性的，完全平行于
\cite[推论 7.1.8]{ABC+24a} 可以断定
\[
M^{\temp}\simeq
\Fun_{\QCoh((\LS^{\mer}_{\check G,x_0})^{\wedge\mathrm{nilp\text{-}mon}})}
\Bigl(\QCoh((\check{\mathfrak{n}}/\check B)_{x_0}),
\QCoh((\check{\mathfrak{n}}/\check B)_{x_0})
\otimes_{\Aff^{\spec}_{\check G,x_0}}M\Bigr);
\]
从而在 Iwahori–Hecke 温的 $\Aff^{\spec}_{G,x_0}$-模范畴的全子范畴上，对应
\[
M\longmapsto \QCoh((\check{\mathfrak{n}}/\check B)_{x_0})
\otimes_{\Aff^{\spec}_{\check G,x_0}}M
\]
是保守的。

下面是我们证明命题~\ref{prop:3.4.2.1} 的策略。它类似于 \cite[第 7.2 节至第 7.4
节]{ABC+24a} 中的论证。首先，我们将构造交换图
\[
\begin{tikzcd}[column sep=large,row sep=large,cells={font=\small}]
\Whit^{*}(\Fl_G)_{x_0}\otimes_{\Aff_{G,x_0}}\IKL(G)_{\crit,x_0}
\arrow[r] \arrow[d]
& \QCoh(\check{\mathfrak{n}}/\check B)_{x_0}
\otimes_{\Aff^{\spec}_{G,x_0}}
\IndCoh^{*}\bigl(\Op^{\mer}_{\check G}
\times_{\LS^{\mer}_{\check G}}\check{\mathfrak{n}}/\check B\bigr)_{x_0}
\arrow[d]\\
\IndCoh(\Op^{\mer}_{\check G,x_0}) \arrow[r,equal] &
\IndCoh(\Op^{\mer}_{\check G,x_0}).
\end{tikzcd}
\]
然后说明两条竖箭头都是全忠实嵌入且本质像相同，这将表明顶部箭头
\begin{equation}
\Whit^{*}(\Fl_G)_{x_0}\otimes_{\Aff_{G,x_0}}\IKL(G)_{\crit,x_0}
\longrightarrow
\QCoh(\check{\mathfrak{n}}/\check B)_{x_0}
\otimes_{\Aff^{\spec}_{G,x_0}}\IKL(G)_{\crit,x_0}
\label{eq:3.4.5}
\end{equation}
\begin{equation}
\xrightarrow{\ \mathrm{Id}\otimes \IFLE_{x_0}\ }
\QCoh(\check{\mathfrak{n}}/\check B)_{x_0}
\otimes_{\Aff^{\spec}_{G,x_0}}
\IndCoh^{*}\bigl(\Op^{\mer}_{\check G}
\times_{\LS^{\mer}_{\check G}}\check{\mathfrak{n}}/\check B\bigr)_{x_0}
\label{eq:3.4.6}
\end{equation}
是等价。由于该复合的第一个态射是等价，这将推出复合的第二个态射也是等价。
\srcpage{113}

最后，我们将说明 $\IKL(G)_{\crit,x_0}$ 与
\[
\IndCoh^{*}\bigl(\Op^{\mer}_{\check G}
\times_{\LS^{\mer}_{\check G}}\check{\mathfrak{n}}/\check B\bigr)_{x_0}
\simeq
\IndCoh\bigl(\Op^{\mer}_{\check G,x_0}
\times_{\LS^{\mer}_{\check G,x_0}}(\check{\mathfrak{n}}/\check B)_{x_0}\bigr)
\]
都是 Iwahori–Hecke 温的，从而可以断定 $\IFLE_{x_0}$ 是等价。

\subsection{IFLE 是逐点等价}
\label{sec:ifle-pointwise}

本小节逐一验证为证明命题~\ref{prop:3.4.2.1} 所需的各种断言。

\begin{proposition}\label{prop:3.4.4.1}
存在交换图
\[
\begin{tikzcd}[column sep=large,row sep=large,cells={font=\small}]
\Whit^{*}(\Fl_G)_{x_0}\otimes_{\Aff_{G,x_0}}\IKL(G)_{\crit,x_0}
\arrow[r] \arrow[d]
& \QCoh(\check{\mathfrak{n}}/\check B)_{x_0}
\otimes_{\Aff^{\spec}_{G,x_0}}
\IndCoh^{*}\bigl(\Op^{\mer}_{\check G}
\times_{\LS^{\mer}_{\check G}}\check{\mathfrak{n}}/\check B\bigr)_{x_0}
\arrow[d]\\
\IndCoh(\Op^{\mer}_{\check G,x_0}) \arrow[r,equal] &
\IndCoh(\Op^{\mer}_{\check G,x_0}).
\end{tikzcd}
\]
\end{proposition}

\begin{proof}
我们先断言 $\FactModCat_{x_0}$ 中有交换图
\[
\begin{tikzcd}[column sep=large,row sep=large,cells={font=\small}]
\IKL(G)_{\crit} \arrow[r] \arrow[d]
& \IndCoh^{*}\bigl(\Op^{\mer}_{\check G}
\times_{\LS^{\mer}_{\check G}}\check{\mathfrak{n}}/\check B\bigr) \arrow[d]\\
\Whit^{!}(\Fl_G)\otimes \IndCoh^{*}(\Op^{\mer}_{\check G,x_0}) \arrow[r] &
\QCoh(\check{\mathfrak{n}}/\check B)\otimes \IndCoh^{*}(\Op^{\mer}_{\check G,x_0}),
\end{tikzcd}
\]
\srcpage{114}
其中右侧竖箭头由沿
\[
\Op^{\mer}_{\check G}\times_{\LS^{\mer}_{\check G}}\check{\mathfrak{n}}/\check B
\longrightarrow \Op^{\mer}_{\check G}\times \check{\mathfrak{n}}/\check B
\]
的 $*$-推前给出。

至于左侧竖箭头，考虑如下范式：若 $M$ 是带 $D\text{-}\mathrm{mod}_{1/2}(LG)$ 作用的范畴，
则有映射
\[
D\text{-}\mathrm{mod}_{1/2}(\Fl^G)\otimes_{\Aff_G}M^{I}\longrightarrow M;
\]
取 Whittaker 余不变量得到
\[
\Whit^{*}(\Fl_G)\otimes_{\Aff_G}M^{I}\longrightarrow \Whit^{*}(M),
\]
再由对偶得到
\[
M^{I}\longrightarrow \Whit^{!}(\Fl_G)\otimes \Whit^{*}(M).
\]
左侧竖箭头就是取 $M=\hat{\mathfrak{g}}\text{-}\mathrm{mod}_{\mathrm{crit}}$ 时的这个函子。

上述交换图可以按 \cite[附录 E.10]{ABC+24a} 的构造方式逐字升级到对的情形而得到。
\end{proof}

于是可以引用 AB 与 IFLE 和 B 的相容性作 curry，得到
\[
\begin{tikzcd}[column sep=large,row sep=large,cells={font=\small}]
\Whit^{*}(\Fl_G)\otimes_{\Aff_G}\IKL(G)_{\crit} \arrow[r] \arrow[d]
& \QCoh(\check{\mathfrak{n}}/\check B)\otimes_{\Aff^{\spec}_G}
\IndCoh^{*}\bigl(\Op^{\mer}_{\check G}
\times_{\LS^{\mer}_{\check G}}\check{\mathfrak{n}}/\check B\bigr) \arrow[d]\\
\IndCoh^{*}(\Op^{\mer}_{\check G}) \arrow[r,equal] &
\IndCoh^{*}(\Op^{\mer}_{\check G}).
\end{tikzcd}
\]
\srcpage{115}
所要的图就是上述图在 $x_0$ 处的特殊化。

\begin{proposition}\label{prop:3.4.4.2}
右侧竖箭头
\[
\QCoh(\check{\mathfrak{n}}/\check B)_{x_0}
\otimes_{\Aff^{\spec}_{G,x_0}}
\IndCoh\bigl(\Op^{\mer}_{\check G,x_0}
\times_{\LS^{\mer}_{\check G,x_0}}(\check{\mathfrak{n}}/\check B)_{x_0}\bigr)
\longrightarrow \IndCoh(\Op^{\mer}_{\check G,x_0})
\]
是全忠实的，其本质像是
\[
\IndCoh(\Op^{\mer}_{\check G,x_0})^{\Op^{\mathrm{nilp\text{-}mon}}_{\check G,x_0}}
\subset \IndCoh(\Op^{\mer}_{\check G,x_0}),
\]
即 $\IndCoh(\Op^{\mer}_{\check G,x_0})$ 中在
\[
\Op^{\mathrm{nilp\text{-}mon}}_{\check G,x_0}
=\Op^{\mer}_{\check G,x_0}
\times_{\LS^{\mer}_{\check G,x_0}}
\LS^{\mathrm{nilp\text{-}mon}}_{\check G,x_0}
\subset \Op^{\mer}_{\check G,x_0}
\]
上集论支撑的对象构成的全子范畴。
\end{proposition}

\begin{proof}
$\QCoh(\check{\mathfrak{n}}/\check B)_{x_0}$ 与 $\Aff^{\spec}_{G,x_0}$ 都和它们的球面
对应物一样是刚性的，故可基本照搬 \cite[命题 3.6.5]{ABC+24a} 证明中的论证：函子
\begin{equation}
\IndCoh\bigl(\Op^{\mer}_{\check G,x_0}
\times_{\LS^{\mer}_{\check G,x_0}}(\check{\mathfrak{n}}/\check B)_{x_0}\bigr)
\simeq \QCoh(\check{\mathfrak{n}}/\check B)_{x_0}
\otimes_{\QCoh(\check{\mathfrak{n}}/\check B)_{x_0}}
\IndCoh\bigl(\Op^{\mer}_{\check G,x_0}
\times_{\LS^{\mer}_{\check G,x_0}}(\check{\mathfrak{n}}/\check B)_{x_0}\bigr)
\label{eq:3.4.7}
\end{equation}
\begin{equation}
\to \QCoh(\check{\mathfrak{n}}/\check B)_{x_0}
\otimes_{\Aff^{\spec}_{G,x_0}}
\IndCoh\bigl(\Op^{\mer}_{\check G,x_0}
\times_{\LS^{\mer}_{\check G,x_0}}(\check{\mathfrak{n}}/\check B)_{x_0}\bigr)
\label{eq:3.4.8}
\end{equation}
\srcpage{116}
有连续且单道的右伴随。推前
\[
\IndCoh\bigl(\Op^{\mer}_{\check G,x_0}
\times_{\LS^{\mer}_{\check G,x_0}}(\check{\mathfrak{n}}/\check B)_{x_0}\bigr)
\longrightarrow
\IndCoh(\Op^{\mer}_{\check G,x_0})^{\Op^{\mathrm{nilp\text{-}mon}}_{\check G,x_0}}
\]
也是如此，因为
\[
\check{\mathfrak{n}}/\check B\longrightarrow \LS^{\mathrm{nilp\text{-}mon}}_{\check G,x_0}
\]
是 Springer 消解，从而是真且满的。函子
\[
\QCoh(\check{\mathfrak{n}}/\check B)_{x_0}
\otimes_{\Aff^{\spec}_{G,x_0}}
\IndCoh\bigl(\Op^{\mer}_{\check G,x_0}
\times_{\LS^{\mer}_{\check G,x_0}}(\check{\mathfrak{n}}/\check B)_{x_0}\bigr)
\longrightarrow \IndCoh(\Op^{\mer}_{\check G,x_0})
\]
诱导两个单子之间的态射；要说明它是单子的同构，只需说明它在自函子层面是同构。
然而两个单子都由代数对象
\[
\Bigl(\mathrm{pr}_1\colon (\check{\mathfrak{n}}/\check B)_{x_0}
\times_{\LS^{\mer}_{\check G,x_0}}(\check{\mathfrak{n}}/\check B)_{x_0}
\to (\check{\mathfrak{n}}/\check B)_{x_0}\Bigr)^{!}
\mathcal{O}_{(\check{\mathfrak{n}}/\check B)_{x_0}}
\]
（它是 $\Aff^{\spec}_{\check G,x_0}$ 的代数对象）在
$\IndCoh\bigl(\Op^{\mer}_{\check G,x_0}
\times_{\LS^{\mer}_{\check G,x_0}}(\check{\mathfrak{n}}/\check B)_{x_0}\bigr)$
上的作用给出。
\end{proof}

\begin{proposition}\label{prop:3.4.4.3}
左侧竖箭头
\[
\Whit^{*}(\Fl_G)_{x_0}\otimes_{\Aff_{G,x_0}}\IKL(G)_{\crit,x_0}
\longrightarrow \IndCoh(\Op^{\mer}_{\check G,x_0})
\]
是全忠实的，其本质像是
\[
\IndCoh(\Op^{\mer}_{\check G,x_0})^{\Op^{\mathrm{nilp\text{-}mon}}_{\check G,x_0}}
\subset \IndCoh(\Op^{\mer}_{\check G,x_0}).
\]
\end{proposition}

\begin{proof}
与仿射 Grassmann 叠一样，仿射旗簇 $\Fl_{\check G,x_0}$ 也是 ind-真的。因此全忠实性的
证明类似于 \cite[第 7.3 节]{ABC+24a} 中的论证：与球面情形一样，对临界层级上带环路群
$LG$ 作用的范畴 $C$，
\[
C^{\mathrm{Aff\text{-}gen}}
=D\text{-}\mathrm{mod}_{1/2}(\Fl_{G,x_0})
\otimes_{\Aff_{G,x_0}}C^{I}\longrightarrow C
\]
是全忠实的且有连续右伴随。但这推出
\[
\Whit^{*}(C^{\mathrm{Aff\text{-}gen}})
=\Whit^{*}(\Fl_G)_{x_0}\otimes_{\Aff_{G,x_0}}C^{I}
\longrightarrow \Whit^{*}(C)
\]
是全忠实的。取 $C=\hat{\mathfrak{g}}\text{-}\mathrm{mod}_{\mathrm{crit}}$，即得
\[
\Whit^{*}(\Fl_G)_{x_0}\otimes_{\Aff_{G,x_0}}\IKL(G)_{\crit,x_0}
\longrightarrow
\Whit^{*}(\hat{\mathfrak{g}}\text{-}\mathrm{mod}_{\mathrm{crit}})
\simeq \IndCoh(\Op^{\mer}_{\check G,x_0})
\]
的全忠实性，此即所需。

此外，函子
\[
\IKL(G)_{\crit,x_0}\xrightarrow{\ c\,\mapsto\, 1^{\Aff}_G\otimes c\ }
\Whit^{*}(\Fl_G)_{x_0}\otimes_{\Aff_{G,x_0}}\IKL(G)_{\crit,x_0}
\]
生成目标；因此要说明
\[
\Whit^{*}(\Fl_G)_{x_0}\otimes_{\Aff_{G,x_0}}\IKL(G)_{\crit,x_0}
\longrightarrow \IndCoh(\Op^{\mer}_{\check G,x_0})
\]
的本质像落在
$\IndCoh(\Op^{\mer}_{\check G,x_0})^{\Op^{\mathrm{nilp\text{-}mon}}_{\check G,x_0}}$
中，只需说明复合
\[
\IKL(G)_{\crit,x_0}\longrightarrow
\Whit^{*}(\Fl_G)_{x_0}\otimes_{\Aff_{G,x_0}}\IKL(G)_{\crit,x_0}
\longrightarrow \IndCoh(\Op^{\mer}_{\check G,x_0})
\]
的本质像落在其中。然而上述复合同构于复合
\srcpage{118}
\[
\IKL(G)_{\crit,x_0}\xrightarrow{\ \IFLE_{x_0}\ }
\IndCoh\bigl(\Op^{\mer}_{\check G,x_0}
\times_{\LS^{\mer}_{\check G,x_0}}(\check{\mathfrak{n}}/\check B)_{x_0}\bigr)
\longrightarrow \IndCoh(\Op^{\mer}_{\check G,x_0}),
\]
后者的本质像显然落在
$\IndCoh(\Op^{\mer}_{\check G,x_0})^{\Op^{\mathrm{nilp\text{-}mon}}_{\check G,x_0}}$ 中。

现在剩下要说明
\[
\Whit^{*}(\Fl_G)_{x_0}\otimes_{\Aff_{G,x_0}}\IKL(G)_{\crit,x_0}
\longrightarrow
\IndCoh(\Op^{\mer}_{\check G,x_0})^{\Op^{\mathrm{nilp\text{-}mon}}_{\check G,x_0}}
\]
本质满。由于该函子全忠实，只需说明它生成目标；而这又只需说明复合
\[
\IKL(G)_{\crit,x_0}\longrightarrow
\Whit^{*}(\Fl_G)_{x_0}\otimes_{\Aff_{G,x_0}}\IKL(G)_{\crit,x_0}
\longrightarrow
\IndCoh(\Op^{\mer}_{\check G,x_0})^{\Op^{\mathrm{nilp\text{-}mon}}_{\check G,x_0}}
\]
生成目标。

由于有
\[
\bigl(\Op^{\mathrm{nilp\text{-}mon}}_{\check G,x_0}\bigr)_{\mathrm{red}}
=\bigcup_{\check\lambda\in \check\lambda^{+}}
\Op^{\check\lambda\text{-}\mathrm{nilp}}_{\check G,x_0},
\]
诸对象
\[
\mathcal{O}_{\Op^{\check\lambda\text{-}\mathrm{nilp}}_{\check G,x_0}}
\in
\IndCoh(\Op^{\mer}_{\check G,x_0})^{\Op^{\mathrm{nilp\text{-}mon}}_{\check G,x_0}}
\]
紧生成
$\IndCoh(\Op^{\mer}_{\check G,x_0})^{\Op^{\mathrm{nilp\text{-}mon}}_{\check G,x_0}}$。

我们断言：Verma 模
\[
M^{\check\lambda}=\mathrm{ind}(\hat{\mathfrak{g}}^{\mathrm{crit}},I)_I\,
k^{\check\lambda}\in \IKL(G)_{\crit,x_0}
\]
在函子
$\IKL(G)_{\crit,x_0}\to
\IndCoh(\Op^{\mer}_{\check G,x_0})^{\Op^{\mathrm{nilp\text{-}mon}}_{\check G,x_0}}$
下被送到
$\mathcal{O}_{\Op^{\check\lambda\text{-}\mathrm{nilp}}_{\check G,x_0}}$。

为看出这一点，我们将改装 \cite{FG10} 中证明 FLE 把 Weyl 模 $V^{\check\lambda}$ 送到
$\mathcal{O}_{\Op^{\check\lambda\text{-}\mathrm{reg}}_{\check G,x_0}}$ 的论证：
\srcpage{119}
一方面有同构
\[
\mathcal{O}_{\Op^{\check\lambda\text{-}\mathrm{nilp}}_{\check G,x_0}}
\simeq \End(M^{\check\lambda}),
\]
而展开定义可得
\[
\End(M^{\check\lambda})
=\Hom\bigl(\mathrm{ind}(\hat{\mathfrak{g}}^{\mathrm{crit}},I)_I k^{\check\lambda},
M^{\check\lambda}\bigr)
\simeq \Hom_I(k^{\check\lambda},M^{\check\lambda})
=:(M^{\check\lambda})^{I^0_{\check\lambda}}.
\]
对任何 $M\in \hat{\mathfrak{g}}\text{-}\mathrm{mod}_{\mathrm{crit}}$ 都有映射
\[
(M)^{I^0_{\check\lambda}}\longrightarrow M^{\hat{\mathfrak{n}}_+}
\longrightarrow H^{\infty}_2\bigl(n((t)),n[[t]],M\otimes \Psi_0\bigr),
\]
这正是 \cite[引理 3]{FG10} 的内容。取 $M=M^{\check\lambda}$ 即得映射
\[
\mathcal{O}_{\Op^{\check\lambda\text{-}\mathrm{nilp}}_{\check G,x_0}}
\longrightarrow
H^{\infty}_2\bigl(n((t)),n[[t]],M^{\check\lambda}\otimes \Psi_0\bigr),
\]
而右端可等同于 $M^{\check\lambda}$ 的所要像。因此只需说明该映射是同构。

与 $\check\lambda$-正则情形类似，上述映射具有
$N^{*}_{\Op^{\check\lambda\text{-}\mathrm{nilp}}_{\check G,x_0}/
\Op^{\mer}_{\check G,x_0}}$-模结构，且
$\mathcal{O}_{\Op^{\check\lambda\text{-}\mathrm{nilp}}_{\check G,x_0}}$ 作为
$N^{*}_{\Op^{\check\lambda\text{-}\mathrm{nilp}}_{\check G,x_0}/
\Op^{\mer}_{\check G,x_0}}$-模是不可约的。因此该映射必为单射，只需按
\cite[第 5 节]{FG10} 计算并比较两边的特征。

先计算
$\mathcal{O}_{\Op^{\check\lambda\text{-}\mathrm{nilp}}_{\check G,x_0}}$ 的特征。记
\[
J=\{\alpha_i\in \Delta\mid \langle \alpha_i,\check\lambda+\check\rho\rangle=0\},
\]
并记 $\check P_J$ 为 $\check G$ 中对应于 $J$ 的抛物子群。注意
$\Op^{\check\lambda\text{-}\mathrm{nilp}}_{\check G,x_0}$ 是形如
\srcpage{120}
\[
\Bigl\{\partial_t+\sum_{\Delta} t^{\langle \alpha_i,\check\lambda\rangle}f_i
+v(t)+\frac{w}{t}\ \Big|\ v(t)\in \check{\mathfrak{b}}[[t]],\ w\in O\Bigr\}
\]
的算子之集，模去 $\check N[[t]]$ 的作用，其中
\[
O=\Bigl\{w\in \check{\mathfrak{b}}\ \Big|\
w+\sum_{J}f_i \bmod \check{\mathfrak{n}}_J\in N_{\check{\mathfrak{m}}_J}\Bigr\}.
\]
然而这与
\begin{align*}
&\Bigl\{\partial_t+\frac{1}{t}\Bigl(\sum^{\Delta}f^i-(\check\lambda+\check\rho)\Bigr)
+v(t)+\frac{w}{t}\ \Big|\
v(t)\in (\check\lambda+\check\rho)(t)\check{\mathfrak{b}}[[t]](\check\lambda+\check\rho)(t)^{-1},\\
&\qquad\qquad\qquad\qquad\qquad
w\in (\check\lambda+\check\rho)(t)O(\check\lambda+\check\rho)(t)^{-1}\Bigr\}
\end{align*}
相同，模去 $(\check\lambda+\check\rho)(t)\check N[[t]](\check\lambda+\check\rho)(t)^{-1}$
的自由作用。
\yizhu{原文第一式中 $f_i$ 前的指数为 $\langle\alpha_i,\check\lambda\rangle$，
第二式中相应部分写作 $\frac{1}{t}\bigl(\sum^{\Delta}f^i-(\check\lambda+\check\rho)\bigr)$，
正文称两式「相同」。令 $A=\partial_t+\frac{1}{t}\bigl(\sum^{\Delta}f^i-
(\check\lambda+\check\rho)\bigr)$，$g=(\check\lambda+\check\rho)(t)$，则
$\mathrm{Ad}(g)A=\partial_t+\sum_{\Delta}t^{\langle\alpha_i,\check\lambda+
\check\rho\rangle-1}f_i$。故第一式的指数应为
$\langle\alpha_i,\check\lambda+\check\rho\rangle-1$；原文如此，恐系笔误，
兹照排并记于此。}{译者按（验算所得，无外部依据）}

考虑 $O\cap \check{\mathfrak{m}}_J$，它带有群 $\check N\cap \check M_J$ 的自由传递
作用。有群的满同态
\[
\phi\colon (\check\lambda+\check\rho)(t)\check N[[t]](\check\lambda+\check\rho)(t)^{-1}
\longrightarrow \check N\cap \check M_J,
\]
它取出常数项；并有满态射
\begin{align}
&\Bigl\{\partial_t+\frac{1}{t}\Bigl(\sum^{\Delta}f^i-(\check\lambda+\check\rho)\Bigr)
+v(t)+\frac{w}{t}\ \Big|\
v(t)\in (\check\lambda+\check\rho)(t)\check{\mathfrak{b}}[[t]](\check\lambda+\check\rho)(t)^{-1},\notag\\
&\qquad\qquad\qquad
w\in (\check\lambda+\check\rho)(t)O(\check\lambda+\check\rho)(t)^{-1}\Bigr\}
\label{eq:3.4.9}\\
&\xrightarrow{\ \pi\ } O\cap \check{\mathfrak{m}}_J,
\label{eq:3.4.10}
\end{align}
它取出 $w$ 的常数项。$\pi$ 关于 $\phi$ 是等变的，于是
\[
\Op^{\check\lambda\text{-}\mathrm{nilp}}_{\check G,x_0}
\simeq \pi^{-1}(0)/\ker(\phi),
\]
\srcpage{121}
从而
\[
\operatorname{char}\bigl(\mathcal{O}_{\Op^{\check\lambda\text{-}\mathrm{nilp}}
_{\check G,x_0}}\bigr)
=\operatorname{char}\bigl(\Fun(\pi^{-1}(0))\bigr)/
\operatorname{char}\bigl(\Fun(\ker(\phi))\bigr).
\]
而 $\pi^{-1}(0)$ 是
\begin{align*}
&\Bigl\{\partial_t+\frac{1}{t}\Bigl(\sum^{\Delta}f^i-(\check\lambda+\check\rho)\Bigr)
+v(t)+\frac{w}{t}\ \Big|\
v(t)\in (\check\lambda+\check\rho)(t)\check{\mathfrak{b}}[[t]](\check\lambda+\check\rho)(t)^{-1},\\
&\qquad\qquad\qquad\qquad\qquad
w\in (\check\lambda+\check\rho)(t)\check{\mathfrak{n}}^J(\check\lambda+\check\rho)(t)^{-1}\Bigr\},
\end{align*}
故
\[
\operatorname{char}\bigl(\Fun(\pi^{-1}(0))\bigr)
=\prod_{n>0}(1-q)^{-l}\cdot
\Bigl(\prod_{\alpha\in \Phi^{+}_J}\prod_{n\geq 0}(1-q^{n+1})^{-1}\Bigr)\cdot
\Bigl(\prod_{\alpha\in \Phi^{+}\setminus \Phi^{+}_J}
\prod_{n\geq 0}(1-q^{n+\langle \alpha,\check\lambda+\check\rho\rangle})^{-1}\Bigr).
\]
另一方面，
\[
\ker(\phi)=(\check\lambda+\check\rho)(t)\check N[[t]](\check\lambda+\check\rho)(t)^{-1}
\cap \exp(t\check{\mathfrak{n}}[[t]]),
\]
故
\[
\operatorname{char}\bigl(\Fun(\ker(\phi))\bigr)
=\Bigl(\prod_{\alpha\in \Phi^{+}_J}\prod_{n\geq 0}(1-q^{n+1})^{-1}\Bigr)\cdot
\Bigl(\prod_{\alpha\in \Phi^{+}\setminus \Phi^{+}_J}
\prod_{n\geq 0}(1-q^{n+\langle \alpha,\check\lambda+\check\rho\rangle})^{-1}\Bigr).
\]
因此
\[
\operatorname{char}\bigl(\mathcal{O}_{\Op^{\check\lambda\text{-}\mathrm{nilp}}
_{\check G,x_0}}\bigr)=\prod_{n>0}(1-q)^{-l}.
\]

最后，为计算
$H^{\infty}_2\bigl(n((t)),n[[t]],M^{\check\lambda}\otimes \Psi_0\bigr)$ 的特征，
我们基本逐字照搬 \cite[第 5.2 节]{FG10} 的做法，只把有限 Weyl 模 $V^{\lambda}$ 换成有限
Verma 模 $M^{\lambda}$。于是得到
\[
\operatorname{char}\bigl(H^{\infty}_2\bigl(n((t)),n[[t]],
M^{\check\lambda}\otimes \Psi_0\bigr)\bigr)=\prod_{n>0}(1-q)^{-l},
\]
此即所求。
\end{proof}
\srcpage{122}

\begin{proposition}\label{prop:3.4.4.4}
$\IndCoh\bigl(\Op^{\mer}_{\check G,x_0}
\times_{\LS^{\mer}_{\check G,x_0}}(\check{\mathfrak{n}}/\check B)_{x_0}\bigr)$
是 Iwahori–Hecke 温的。
\end{proposition}

\begin{proof}
由于平凡的理由，有同构
\[
\Op^{\mer}_{\check G,x_0}
\times_{\LS^{\mer}_{\check G,x_0}}(\check{\mathfrak{n}}/\check B)_{x_0}
\simeq
(\Op^{\mer}_{\check G,x_0})^{\wedge\mathrm{nilp\text{-}mon}}
\times_{(\LS^{\mer}_{\check G,x_0})^{\wedge\mathrm{nilp\text{-}mon}}}
(\check{\mathfrak{n}}/\check B)_{x_0}.
\]
由于 $(\LS^{\mer}_{\check G,x_0})^{\wedge\mathrm{nilp\text{-}mon}}$ 与它的球面对应物
$(\LS^{\mer}_{\check G,x})^{\wedge\mathrm{mon\text{-}free}}$ 一样是可通行的，作为
$\Aff^{\spec}_{G,x_0}$-模范畴有
\begin{equation}
\IndCoh\bigl((\Op^{\mer}_{\check G,x_0})^{\wedge\mathrm{nilp\text{-}mon}}
\times_{(\LS^{\mer}_{\check G,x_0})^{\wedge\mathrm{nilp\text{-}mon}}}
(\check{\mathfrak{n}}/\check B)_{x_0}\bigr)
\label{eq:3.4.11}
\end{equation}
\begin{equation}
\simeq
\IndCoh\bigl((\Op^{\mer}_{\check G,x_0})^{\wedge\mathrm{nilp\text{-}mon}}\bigr)
\otimes_{\QCoh((\LS^{\mer}_{\check G,x_0})^{\wedge\mathrm{nilp\text{-}mon}})}
\QCoh(\check{\mathfrak{n}}/\check B)_{x_0}.
\label{eq:3.4.12}
\end{equation}
然而 $\Aff^{\spec}_{G,x_0}$ 在右端上的作用来自它在
$\QCoh(\check{\mathfrak{n}}/\check B)_{x_0}$ 上的作用，而后者经
$\Aff^{\spec,\temp}_{\check G,x_0}$ 分解。
\end{proof}

最后，我们说明 $\IKL(G)_{\crit,x_0}$ 是 Iwahori–Hecke 温的。

\begin{remark}\label{rem:3.4.4.5}
此处的论证与 \cite[命题 7.2.6]{ABC+24a} 的相应结果在 \cite[第 7.4 节]{ABC+24a} 中给出的
证明有显著差别。原因是 $\Aff^{\spec}_{G,x_0}$ 在 $\IKL(G)_{\crit,x_0}$ 上的作用不是
t-正合的。我们改用下面这个较弱的性质。
\end{remark}

\begin{lemma}\label{lem:3.4.4.6}
若
\[
F\in \Aff^{\spec,\leq -\infty}_{\check G,x_0},
\qquad
c\in \IKL(G)^{c}_{\crit,x_0}\subset \IKL(G)^{b}_{\crit,x_0},
\]
则
\[
F\star c\in \IKL(G)^{\leq -\infty}_{\crit,x_0}.
\]
\end{lemma}

\begin{proof}
由于 $\IKL(G)_{\crit,x_0}$ 由 Verma 模生成，可设 $c=M^{\check\lambda}$。函子
\[
\Gamma\colon D\text{-}\mathrm{mod}_{\mathrm{crit}}(\Fl_G)\longrightarrow
\IKL(G)_{\crit,x_0}
\]
在 \cite{FG09} 中已有研究。对对象 $F\in \Aff_G$ 有
\[
\Gamma(F)=F\star M^{0}.
\]
用类似方法可以研究扭版本
\[
\Gamma^{\check\lambda}\colon
D\text{-}\mathrm{mod}_{\mathrm{crit},\check\lambda}(\Fl_G)\longrightarrow
\IKL(G)_{\crit,x_0},
\]
它把 $F$ 送到
\[
\Gamma^{\check\lambda}(F)=F\star M^{\check\lambda}.
\]
于是断言由每个 $\Gamma^{\check\lambda}$ 的有界上同调振幅推出。
\end{proof}
\srcpage{123}

\begin{proposition}\label{prop:3.4.4.7}
$\IKL(G)_{\crit,x_0}$ 是 Iwahori–Hecke 温的。
\end{proposition}

\begin{proof}
要说明
\[
\Temp\colon \IKL(G)_{\crit,x_0}\longrightarrow
\IKL(G)^{\temp}_{\crit,x_0}
\]
是等价。由于它是余局部化，只需说明它是保守的。注意函子
\[
\IFLE^{\temp}\colon
\IKL(G)^{\temp}_{\crit,x_0}\longrightarrow
\IndCoh\bigl(\Op^{\mer}_{\check G,x_0}
\times_{\LS^{\mer}_{\check G,x_0}}(\check{\mathfrak{n}}/\check B)_{x_0}\bigr)^{\temp},
\]
它与
\[
\IndCoh\bigl(\Op^{\mer}_{\check G,x_0}
\times_{\LS^{\mer}_{\check G,x_0}}(\check{\mathfrak{n}}/\check B)_{x_0}\bigr)
\]
等同。
是等价。因此，由于 IFLE 保持紧对象，$\Temp$ 也保持紧对象。于是只需说明 $\Temp$ 在紧对象
上是保守的。然而 $\IKL(G)_{\crit,x_0}$ 中所有紧对象都是有界的，并且利用
\[
\operatorname{cone}\bigl(\Temp(1_{\Aff^{\spec}_{G,x_0}})
\to 1_{\Aff^{\spec}_{G,x_0}}\bigr)
\in \Aff^{\spec,\leq -\infty}_{\check G,x_0},
\]
\srcpage{124}
可以断定余单位映射的锥
\[
\operatorname{cone}\bigl(\Temp(c)\to c\bigr)
\simeq
\operatorname{cone}\bigl(\Temp(1_{\Aff^{\spec}_{G,x_0}})
\to 1_{\Aff^{\spec}_{G,x_0}}\bigr)\star c
\in \IKL(G)^{\leq -\infty}_{\crit,x_0}.
\]
于是若 $\Temp(c)=0$，则
$c\in \IKL(G)^{\leq -\infty}_{\crit,x_0}$，而 $c$ 又有界，故 $c=0$。这就说明 $\Temp$
在紧对象上保守，此即所求。
\end{proof}
